\section{Discussion}
\label{sec:discussion}

\subsection{Why Each Branch Helps}

% TODO: Fill in after ablation results.
% Expected argument structure:
% - Graph catches structural patterns (e.g., tainted data flowing to sink)
% - Image catches global layout patterns (e.g., deeply nested control flow)
% - LLM catches semantic context (e.g., unsafe API usage patterns)

\subsection{When VulGCL Fails}

% TODO: Analyze false positives and false negatives after experiments.
% Expected: obfuscated code, very short functions, unusual coding styles.

\subsection{Limitations}

\begin{itemize}
  \item VulGCL is currently evaluated on C/C++ (Devign, BigVul). Generalization
    to other languages requires Joern support and retraining.
  \item PDG extraction via Joern fails on syntactically malformed code snippets
    (approx.\ X\% of SO snippets in our dataset).
  \item VulGCL-Lite's quantization causes a modest accuracy drop of
    approximately X\%, which may be unacceptable in high-security contexts.
\end{itemize}

\subsection{Stack Overflow Label Quality}

Raw SO vote scores anticorrelate with security~\cite{chen2019reliable}.
Our composite labeling strategy (static analysis + comment flags + accepted
answer status) addresses this, but static analysis tools have their own false
positive rates. Future work should include manual validation of a sample of SO
labels.
